\chapter{Preparation}

\section{Understanding the matrix-based approach}

{\footnotesize\itshape
The applications of tropical semirings to dynamic dependency graphs, particularly the sections on
register-only longest chains (\ref{prep:register-only-chains})
and alternative graph structures (\ref{prep:alternative-matrix-struct})
are closely adapted from the methodologies established by Chadwick \cite{Chadderz}.}

\subsection{The tropical semiring (max-plus algebra)}

Tropical analysis is a branch of pure mathematics that is the study of the tropical semiring.
There are two variants of the tropical semiring, namely the min-plus and max-plus algebras.
We will use the max-plus convention to apply to \textit{longest} paths in directed acyclic graphs.

 
% Define tropical semiring and operations 

\begin{definition}[Semiring] \label{def:semiring}
  A semiring \cite{Wandsnider2020} is a set $R$ along with two operations, addition (+) and multiplication ($\cdot$), which satisfy the 
  following axioms:

  \begin{enumerate}
    \item Addition is associative and commutative.
    \item Multiplication is associative.
    \item There exists an additive identity, which we usually call $\mathbb{0}$. Additionally, $\mathbb{0}$ annihilates 
      $R$, that is, for all $x \in R$, $\mathbb{0} \cdot x = x \cdot \mathbb{0} = \mathbb{0}$.
    \item There exists a multiplicative identity, which we usually call $\mathbb{1}$.
    \item Multiplication left and right distributes over addition.
  \end{enumerate}
\end{definition}

\begin{definition}[Tropical semiring]
  The tropical semiring \cite{Wandsnider2020}, or max-plus algebra, is the object $(\mathbb{R} \cup \{-\infty\}, \oplus, \odot)$ with operations 
defined as:

$$
\begin{aligned}
  a\oplus b &= \max(a,b) \\
  a\odot b &= a+b
\end{aligned}
$$

\end{definition}

The tropical semiring satisfies the axioms of a semiring from Definition \ref{def:semiring}. It
is not a ring as the max operator does not have an inverse.

% Define matrix multiplication in the context of the tropical semiring

\begin{definition}[Matrix operations in the tropical semiring]
For any $n\times n$ matrices $X = (x_{ij})$ and $Y = (y_{ij})$, the entries of $U = X \oplus Y$ and $V = X \odot Y$
are calculated as \cite{krivulin2012maxplusalgebraapproachmodelling}

$$
u_{ij} = x_{ij} \oplus y_{ij}, \text{ and } v_{ij} = \displaystyle\sum_{k=1}^{n} \! {}_{\oplus} x_{ik} \odot y_{kj}
$$

\end{definition}

The matrix operations $\oplus$ and $\odot$ are associative owing to the commutativity and associativity of the underlying 
addition operator, and that the multiplication operator distributes over addition.

\subsection{Longest paths in Directed Acyclic Graphs}


% Show how this applies to the longest-path problem

As established previously, finding the length of the critical path in a program results in finding the longest path
in its DDG. It is noted that the max-plus version of the tropical semiring is applicable to finding the longest path
in a graph \cite{krivulin2012maxplusalgebraapproachmodelling}.

One algorithm \cite{Chadderz} to find the longest path between nodes $A$ and $B$ in a graph is to recursively take the 
maximum of all path lengths from all neighbours of $A$ to $B$, plus the cost of the edge from $A$ to the chosen
neighbour. We can see that this amounts to a Bellman recursion equation \cite{Bellman:DynamicProgramming}
involving only the max and plus operators: 

$$
v(A,B) = 
\begin{cases}
0 &\text{ if } A = B \\
\max_{\forall C . A \to C} \{\mathrm{weight}(A \to C) + v(C,B))\}, &\text{ otherwise } \\
\end{cases}
$$

We can observe that, if we have an adjacency matrix (where the cells in that matrix represent the cost of the edge
between the pair of nodes in that row/column, and $-\infty$ if there is no edge), performing tropical semiring 
matrix multiplication of this matrix with itself corresponds to one step of our longest-path algorithm.

Formally we consider an adjacency matrix $X^{2} = X \odot X$, and denote its entries by $x_{ij}^{(2)}$. We have that $x_{ij}^{(2)} = p$
if and only if there exists at least one path of length $p$ from node $i$ to node $j$ in the corresponding graph,
consisting of at most two edges, and that path is the longest such path. This extends to any integer exponent $q$
\cite{krivulin2012maxplusalgebraapproachmodelling}.

\begin{theorem}[Matrix exponentiation results in longest paths] \label{thm:mat-exp-longest-paths}

For any integer $q>0$, the matrix
$X^{q}$ has the entry $x_{ij}^{(q)} = p$ only if there is a path of length $p$ with at most $q$ edges from $i$ to $j$,
and that path is the longest such path.

\begin{proof}
  We proceed by induction on $q$, the maximum number of edges in the path.

  \textbf{Base case ($q = 1$):} Consider the matrix $X^1=X$. By the definition of an adjacency matrix, the entry  
  $x_{ij}^{(1)}$ is exactly the weight of the direct edge from node $i$ to node $j$. If no such edge exists, then 
  $x_{ij}^{(1)} = -\infty$. Thus, the base case holds: the entries represent their longest paths consisting of at most 
  one edge.

  \textbf{Inductive step:} Assume the claim holds for $q$. That is the matrix $X^{q}$ has entries $x_{ij}^{q}$ 
  representing the longest path from $i$ to $j$ using at most $q$ edges.
  We show that the claim holds for $q+1$. By definition, $X^{q+1} = X^{q} \odot X$
  The entry $x_{ij}^{(q+1)}$ is calculated as:
  $$x_{ij}^{(q+1)} = \sum_{k=1}^{n} \! {}_{\oplus} (x_{ik}^{(q)}\odot x_{kj})$$
  Translating this to the original operators:
  $$
  x_{ij}^{(q+1)} = \max_{1 \leq k \leq n} (x_{ik}^{(q)} + x_{kj})
  $$
  This lines up perfectly with the dynamic programming approach to finding the longest path. To find a longest path from
  node $i$ to node $j$, it checks every intermediate node $k$. It takes the longest known path from $i$ to $k$ of at most 
  $q$ edges (by the inductive hypothesis), and adds the direct edge cost from $k$ to $j$. By taking the maximum over all 
  intermediate nodes $k$, we find the absolute longest path from $i$ to $j$ of at most $q+1$ edges. Thus the claim holds 
  for all integers $q>0$.
\end{proof}

So, an adjacency matrix raised to the power infinity will converge to the longest distance matrix for that whole graph.
  
\end{theorem}


For example, consider the graph:
\begin{center}
\begin{tikzpicture}[
    node distance=1cm and 2cm, % vertical and horizontal spacing
    every node/.style={font=\Large} % makes letters look like your image
]
    % Place the nodes
    \node (a) {a};
    \node (b) [right=of a] {b};
    \node (c) [below=of b] {c};
    \node (d) [right=of c] {d};

    % Draw the edges (undirected lines) with labels
    \draw [->] (a) -- node[above] {1} (b);
    \draw [->] (a) -- node[below left] {3} (c);
    \draw [->] (b) -- node[right] {3} (c);
    \draw [->] (c) -- node[above] {1} (d);
\end{tikzpicture}
\end{center}

The corresponding adjacency matrix would be:
$$
X = 
\begin{pmatrix}
  0 & 1 & 3 & -\infty \\ 
  -\infty & 0 & 3 & -\infty \\ 
  -\infty & -\infty & 0 & 1 \\
  -\infty & -\infty & -\infty & 0 \\
\end{pmatrix}
$$

So, in the tropical semiring:

$$
\begin{aligned}
  X^{2} &= 
  \begin{pmatrix}
    0 & 1 & 4 & 4 \\ 
    -\infty & 0 & 3 & 4 \\ 
    -\infty & -\infty & 0 & 1 \\ 
    -\infty & -\infty & -\infty & 0 \\
  \end{pmatrix} \\ 
  X^{3} &= 
  \begin{pmatrix}
    0 & 1 & 4 & 5 \\ 
    -\infty & 0 & 3 & 4 \\ 
    -\infty & -\infty & 0 & 1 \\ 
    -\infty & -\infty & -\infty & 0 \\
  \end{pmatrix} \\ 
        & = X^{4} = X^{5} = \dots = X^{\infty}
\end{aligned}
$$

This corresponds to the longest distance matrix between all pairs of nodes.

\subsection{Forming the DDG and register-only longest chains} \label{prep:register-only-chains}

As defined in Section \ref{intro:methodology}, a standard DDG represents 
each dynamically executed instruction as a single node. While this formulation is intuitive for visualising the program 
flow, the structure of the graph is arbitrary, and we would wish to impose a more rigid structure to allow for future
optimisations, as outlined in Section \ref{prep:alternative-matrix-struct}.
To do this, we can encode the transformation of the machine's dependency state from one cycle to the next.
We introduce a new expanded graph structure, which has one node per register per dynamic instruction.

Let us first consider the dependencies for instructions only affecting registers.
Since each instruction has a fixed set of registers which it accesses, it has a fixed dependency structure, and so each 
instruction in the program trace will have a fixed subgraph in the expanded DDG encoding its dependency information.

For example, we could encode the dependencies for the instruction \verb|R0 = ADD R1, R2| as the subgraph:

\begin{center}
\begin{tikzpicture}[
    node distance=0.7cm and 2cm, % vertical and horizontal spacing
    every node/.style={font=\Large} % makes letters look like your image
]
    % Place the nodes
    \node (r0_0) {\verb|R0|$_{0}$};
    \node (r1_0) [below=of r0_0] {\verb|R1|$_{0}$};
    \node (r2_0) [below=of r1_0] {\verb|R2|$_{0}$};
    \node (r0_1) [right=of r0_0] {\verb|R0|$_{1}$};
    \node (r1_1) [below=of r0_1] {\verb|R1|$_{1}$};
    \node (r2_1) [below=of r1_1] {\verb|R2|$_{1}$};

    % Draw the edges
    \draw [->] (r1_0) -- (r0_1);
    \draw [->] (r2_0) -- (r0_1);

    % draw the zero-cost edges
    \draw [->, dashed] (r1_0) -- (r1_1);
    \draw [->, dashed] (r2_0) -- (r2_1);
\end{tikzpicture}

{\tiny dashed edges are 0 cost}
\end{center}

The instruction reads from \verb|R1| and \verb|R2| and uses these values to produce a new value stored in \verb|R0|.
On the left of the subgraph is the register state before the instruction execution; on the right is after. The 
characteristics of the subgraph are as follows:
\begin{itemize}
  \item Since the value \verb|R0|$_1$ is dependent on \verb|R1|$_0$ and \verb|R2|$_0$, there are cost-1 edges in the graph 
    from \verb|R1|$_0$ and \verb|R2|$_0$ to \verb|R0|$_1$, extending the dependency chain from the source registers to the 
    destination by 1 cycle.

  \item The values in \verb|R1| and \verb|R2| are unchanged by the instruction, so their respective states have 0-cost 
    edges between them. This allows dependency chains from the source registers to `skip over' this particular instruction 
    that does not affect them.

  \item Finally, since the value in \verb|R0| is modified, there is no 0-cost edge between the respective states, as the 
    old value is discarded and a new chain starts at the new value.
\end{itemize}


The full DDG will consist of several of these subgraphs, horizontally glued together, with one new set of dependency edges 
for each additional instruction.

The true critical path often involves memory operations, as a value may be copied to memory from a register, then 
potentially millions of instructions later that value is read into a register and used, contributing to a data dependency,
yet we have not accounted for this here. While Chadwick \cite{Chadderz} has suggested approaches to track through-memory
dependencies, this project will not do so, as they require dealing with much larger matrices or significantly complicating 
the structure of the datasets, which fell outside of the project's scope.

Unfortunately, even modest, realistic programs can have trillions of instructions in their 
traces, and so an adjacency matrix for the entire graph becomes unfeasible. However, we can use an alternative 
representation.

\subsection{Alternative graph representation} \label{prep:alternative-matrix-struct}

% Show how we arrange the DDG such that we can have a matrix representing dependencies of an instr.

We can see a particular structure in the DDGs we construct - they are more long than wide. While a 
standard $N \times N$ adjacency matrix (where $N$ is the total number of dynamically executed instructions) could 
map the graph, it is inefficient.
This is because, in the DDG, edges \textit{only} exist between adjacent columns.
Therefore, the full adjacency matrix $X$ is mostly filled with $-\infty$.
Instead, the DDG can be modelled as a sequence of adjacent columns. 
Consider an arbitrary graph with the property of being more long than wide:

\begin{center}
\begin{tikzpicture}[
    node distance=0.7cm and 2cm, % vertical and horizontal spacing
    every node/.style={font=\Large} % makes letters look like your image
]
    % Place the nodes
    \node (a) {a};
    \node (b) [below=of a] {b};
    \node (c) [right=of a] {c};
    \node (d) [below=of c] {d};
    \node (e) [right=of c] {e};
    \node (g) [right=of e] {g};
    \node (f) [below=of e] {f};
    \node (h) [below=of g] {h};

    % Draw the edges (undirected lines) with labels
    \draw [->] (a) -- node[above] {1} (c);
    \draw [->] (a) -- node[below] {4} (d);
    \draw [->] (c) -- node[above] {1} (e);
    \draw [->] (e) -- node[above] {1} (g);
    \draw [->] (f) -- node[above] {2} (g);
    \draw [->] (f) -- node[below] {2} (h);
\end{tikzpicture}
\end{center}

First of all, we modify the graph such that the longest path will always connected to the 
rightmost part of the graph, by adding extra nodes and cost-0 edges, and so also each path is connected to the first 
column similarly.

To utilize column-wise matrices for longest-path calculations, the graph must first be padded. We introduce a set of 
`dummy' nodes to the first and last columns, connected by 0-cost edges. This guarantees that any local path within the 
interior of the graph is extended to touch both the initial and final columns.

\begin{center}
\begin{tikzpicture}[
    node distance=0.7cm and 2cm, % vertical and horizontal spacing
    every node/.style={font=\Large} % makes letters look like your image
]
    % Place the nodes
    \node (a) {a};
    \node (b) [below=of a] {b};
    \node (c) [right=of a] {c};
    \node (d) [below=of c] {d};
    \node (e) [right=of c] {e};
    \node (g) [right=of e] {g};
    \node (f) [below=of e] {f};
    \node (h) [below=of g] {h};

    \node (p1) [above=of a, font=\small] {$p_1$};
    \node (p2) [above=of c, font=\small] {$p_2$};
    \node (p3) [above=of e, font=\small] {$p_3$};
    \node (p4) [above=of g, font=\small] {$p_4$};

    \node (m1) [below=of b, font=\small] {$m_1$};
    \node (m2) [below=of d, font=\small] {$m_2$};
    \node (m3) [below=of f, font=\small] {$m_3$};
    \node (m4) [below=of h, font=\small] {$m_4$};

    % Draw the edges with labels
    \draw [->] (a) -- node[above] {1} (c);
    \draw [->] (a) -- node[below] {4} (d);
    \draw [->] (c) -- node[above] {1} (e);
    \draw [->] (e) -- node[above] {1} (g);
    \draw [->] (f) -- node[above] {2} (g);
    \draw [->] (f) -- node[below] {2} (h);

    % draw the zero-cost edges
    \draw [->, dashed] (p1) -- (p2);
    \draw [->, dashed] (p1) -- (c);
    \draw [->, dashed] (p1) -- (d);

    \draw [->, dashed] (p2) -- (p3);
    \draw [->, dashed] (p2) -- (e);
    \draw [->, dashed] (p2) -- (f);

    \draw [->, dashed] (p3) -- (p4);
    \draw [->, dashed] (p3) -- (g);
    \draw [->, dashed] (p3) -- (h);

    \draw [->, dashed] (a) -- (m2);
    \draw [->, dashed] (b) -- (m2);
    \draw [->, dashed] (m1) -- (m2);

    \draw [->, dashed] (c) -- (m3);
    \draw [->, dashed] (d) -- (m3);
    \draw [->, dashed] (m2) -- (m3);

    \draw [->, dashed] (e) -- (m4);
    \draw [->, dashed] (f) -- (m4);
    \draw [->, dashed] (m3) -- (m4);
\end{tikzpicture}

{\tiny dashed edges are 0 cost}
\end{center}

Now we can represent the graph with one matrix per adjacent-column pair, encoding the connectivity between those two
columns. The connectivity between the first two columns is represented as:

\[
  X_{1,2} = 
  \begin{pmatrix}
  0 & 0 & 0 & -\infty \\
  -\infty & 1 & 4 & 0 \\
  -\infty & -\infty & -\infty & 0 \\
  -\infty & -\infty & -\infty & 0 
  \end{pmatrix}
\]

where the $m$th row and $n$th column of this matrix represent the connection from the $m$th node in the first column of 
the graph to the $n$th node in the second column of the graph.

The remaining two column pairs are represented as:

\[
  \begin{aligned}
    X_{2,3} &= 
    \begin{pmatrix}
    0 & 0 & 0 & -\infty \\
    -\infty & 1 & -\infty & 0 \\
    -\infty & -\infty & -\infty & 0 \\
    -\infty & -\infty & -\infty & 0 
    \end{pmatrix} \\
    X_{3,4} &= 
    \begin{pmatrix}
    0 & 0 & 0 & -\infty \\
    -\infty & 1 & -\infty & 0 \\
    -\infty & 2 & 2 & 0 \\
    -\infty & -\infty & -\infty & 0 
    \end{pmatrix}\\
  \end{aligned}
\]

These matrices are each $4 \times 4$ whereas a full adjacency matrix for the original graph would be $7 \times 7$.
We can see that the column-wise matrices will use less storage than the adjacency matrix (48 integers vs 49 integers),
and this alternative becomes increasingly favourable as the graph gets longer.


In the tropical semiring, computing the matrix product

\[
  X_{3,4} \odot X_{2,3} \odot X_{1,2} = \begin{pmatrix}
  0 & 2 & 2 & 1 \\
  -\infty & 3 & -\infty & 4 \\
  -\infty & -\infty & -\infty & 0 \\
  -\infty & -\infty & -\infty & 0 
\end{pmatrix}
\]
yields a matrix representing the longest path from the $m$th node in the first column the $n$th node in the final 
column of the graph. Because we have ensured every path in the graph connects to the first and last column, the largest 
individual value in this matrix is the length of the longest path in the graph, namely the path of length 4 connecting 
from $a$ in the first column, to $m_4$ in the final column, via $d$ and $m_3$.

The correctness of the column-wise approach for calculating longest paths in graphs is a corollary that follows directly 
from Theorem \ref{thm:mat-exp-longest-paths}.

\begin{corollary}[Correctness of the column-wise formulation]
  In a standard adjacency matrix $X$, computing $X^{q}$ checks every possible intermediate node 
  $k$ in the entire graph to find paths of length $q$, as shown by Theorem \ref{thm:mat-exp-longest-paths}.

  In a column-wise representation of the DDG, when we multiply adjacent column matrices (e.g., $X_{2,3} \odot X_{1,2}$),
  we are simply performing the exact operations for finding the longest path between columns 1 and 3,
  but restricted only to the intermediate nodes $k$ located in column $2$.
  This is valid as edges in the DDG only exist between adjacent columns.

  Because the padding of 0 cost edges ensures every path spans from the first column to the final column, multiplying 
  all consecutive column matrices sequentially computes the maximum path length across the entire graph.
\end{corollary}

This defines the final shape of the data that we will be working on as a mathematical expression consisting of a
series of matrix multiplications. 

We can use the column-wise approach to represent the alternative DDG structure from Section \ref{prep:register-only-chains}.
Because each dynamic instruction has a static set of dependencies, represented by a matrix, and because there are many 
more dynamic instructions in the DDG than the number of static instructions in the program, we will see frequent reuse 
of matrices in the equation, allowing us to find and exploit common patterns, and further reducing the actual memory 
required to store the DDG.

\section{Datasets}

I decided to use execution traces of various coreutils programs as input when testing and benchmarking the analysis tool.
These programs are familiar to developers and display a range of properties that are favourable to several optimisations 
we explore.

\subsection{Dataset listing}

The chosen program executions are:

\begin{itemize}
  \item \verb|cksum < urandom_1MiB.bin|

  the {\color{blue}\textbf{\href{https://www.gnu.org/software/coreutils/manual/html\_node/cksum-invocation.html}{coreutils cksum utility}}} ,
  v8.32 running on AArch64, calculating the CRC of a 1MiB file of random bytes.

  \item \verb|cksum < urandom_32MiB.bin|

  the {\color{blue}\textbf{\href{https://www.gnu.org/software/coreutils/manual/html\_node/cksum-invocation.html}{coreutils cksum utility}}} ,
  v8.32 running on AArch64, calculating the CRC of a 32MiB file of random bytes.

  \item \verb|factor 9696593808888933181586|

  the {\color{blue}\textbf{\href{https://www.gnu.org/software/coreutils/manual/html\_node/factor-invocation.html}{coreutils factor utility}}} ,
  v8.32 running on AArch64, 9696593808888933181586.

  \item \verb|sha1sum < urandom_1MiB.bin|

  the {\color{blue}\textbf{\href{https://www.gnu.org/software/coreutils/manual/html\_node/sha1sum-invocation.html}{coreutils sha1sum utility}}} ,
  v8.32 running on AArch64, computing the SHA1 message digest of a 1MiB file of random bytes.

  \item \verb|sort < urandom_2000_lines.txt|

  the {\color{blue}\textbf{\href{https://www.gnu.org/software/coreutils/manual/html\_node/sort-invocation.html}{coreutils sort utility}}} ,
  v8.32 running on AArch64, sorting a 2000 line text file of hexadecimal digits encoding random bytes.

  \item \verb|shuf --random-source=urandom_1MiB.bin urandom_34953_lines.txt|

  the {\color{blue}\textbf{\href{https://www.gnu.org/software/coreutils/manual/html\_node/shuf-invocation.html}{coreutils shuf utility}}} ,
  v8.32 running on AArch64, shuffling a 34953 line text file of hexadecimal digits encoding random bytes.
\end{itemize}

The source data can be found at \url{https://www.cl.cam.ac.uk/~awc32/cam-only/tropical-ilp-p2p/}. The traces were 
collected by Alexandra Chadwick (\textit{awc32}) using custom tooling.

\subsection{Data properties}

\begin{table}[ht]
\begin{center}
\begin{tabular}{lcccccc}
  \toprule
  \textbf{Dataset} & \textbf{Instructions} & \textbf{Basic Blocks} & \multicolumn{4}{c}{\textbf{Biggest Pattern}\protect\footnotemark} \\
  \cmidrule(l){4-7}
  & (Total) & (Total) & Period\protect\footnotemark & ACI\protect\footnotemark & TI\protect\footnotemark & Coverage\\
  \midrule
  cksum   & 7526926   & 1068940  & 1 & 65535 & 1048560 & 98.1\% \\
  cksum32 & -         & -        & - & - & - & - \\
  factor  & 220632540 & 18074022 & 5 & 20 & 993556 & 27.5\% \\
  sha1sum & 16994993  & 91070    & 1 & 3 & 49155 & 54.0\% \\
  sort    & 10943690  & 1216990  & 17 & 2 & 7536 & 10.5\% \\
  shuf    & 16361747  & 2147424  & 30 & 3 & 16007 & 22.3\% \\
  \bottomrule
\end{tabular}
\end{center}
\caption{Summary of trace analysis on the selected data.}
\label{tbl:trace-analysis}
\end{table}

\footnotetext[1]{Results acquired with an analysis tool based on the loop simplifier in Section IMPLEMENTATION, using a 
window size of 5000 and a history size of 200. The `biggest' pattern in the program is that which covers the largest number 
of basic blocks}
\footnotetext[2]{Period is measured in length of basic blocks and is the length of one iteration of a pattern.}
\footnotetext[3]{Average Consecutive Iterations: The mean number of times a period of a pattern will repeat without interruption.}
\footnotetext[4]{Total Iterations: The total number of times a period of a pattern exists in the program trace.}

As shown in Table \ref{tbl:trace-analysis}, all datasets exhibit a `largest pattern' covering a significant proportion of the trace.
The datasets additionally have other similarly significant patterns not shown in the table.
These patterns can be exploited in two ways:

\begin{enumerate}
  \item Consecutive repetitions of matrices can be quickly reduced down into a single matrix using a quick exponentiation algorithm.
  \item The matrices making up one iteration of the loop can be multiplied together to obtain a singular matrix representing 
    the dependency information for that single iteration. All other iterations of that loop body can then be replaced with 
    the singular result matrix without any extra computation.
\end{enumerate}

\section{Project design}

This section outlines the software engineering methodology adopted during implementation, specification of the 
project requirements, and technical justification for the chosen tools.

\subsection{Software engineering}

For the project, I adopted the iterative spiral methodology. Given the experimental nature of performance optimisation, 
this methodology allowed for a repetitive cycle of planning, followed by implementation, and concluding with an 
evaluation phase where performance bottlenecks were identified to inform the next cycle.

To maintain codebase integrity during high-frequency iterations, code was regularly committed and pushed to GitHub,
using feature branches to implement large changes. Commits from branches were squashed upon merge to provide a clean, 
linear history of major architectural milestones.

\subsection{Project requirements}

% Use a must-should-could pattern for several sections
% MoSCoW

Requirements are listed in Table \ref{tbl:requirements-analysis}
and categorised by their priority, according to the MoSCoW method guidelines \cite{moscow}.

\begin{table}[H]
\centering
\small % Slightly smaller font to fit more content comfortably
\begin{tabularx}{\textwidth}{ll L L}
    \toprule
    \textbf{ID} & \textbf{Category} & \textbf{Requirement} & \textbf{Success Criteria} \\
    \midrule
    NB  & \textit{Must-have} & Naive Baseline Implementation & A valid reference implementation to provide a performance baseline for speedup calculations. \\
    \addlinespace % Adds a tiny bit of breathing room between wrapped rows
    MPF & \textit{Must-have}  & Matrix-Based Path Finding & Correctness verification: ensure the matrix approach yields identical results to the naive baseline. \\
    \addlinespace
    PE  & \textit{Must-have}  & Pattern Exploitation & Demonstrable reduction in redundant computations by identifying and optimising common patterns in the data. \\
    \addlinespace
    MP  & \textit{Should-have} & Multithreaded Parallelism & Achievement of performance scaling on multi-core CPU architectures using an asynchronous sub-problem approach. \\
    \addlinespace
    GA  & \textit{Could-have}  & GPU Acceleration & Successful offloading of matrix multiplications to an accelerator, identifying the point at which throughput overtakes overhead. \\
    \addlinespace
    DC  & \textit{Won't-have}  & Distributed Cluster Compute & Scaling beyond a single node is outside the scope of this project. \\
    \addlinespace
    MT  & \textit{Won't-have}  & Through-memory dependency tracking & Tracking-through memory dependencies is outside the scope of this project. \\
    \bottomrule
\end{tabularx}
\caption{MoSCoW analysis of project requirements and success criteria.}
\label{tbl:requirements-analysis}
\end{table}

\subsection{Project tooling}

The choise of tools was driven by the need for low-level hardware transparency, balanced with architectural flexibility.

C++ was selected as the primary programming language, for three critical reasons:

\begin{enumerate}
  \item It allows for fine-grained memory control, which is essential for performance-sensitive code, allowing for 
    manual cache alignment and minimising overhead from garbage collection pauses.
  \item It provides modern object-oriented paradigms, which allow for the use of design patterns to easily hot-swap 
    different strategies for matrix multiplication and pattern exploitation.
  \item It facilitates the use of SIMD (Single Instruction, Multiple Data) intrinsics for further low-level optimisation.
\end{enumerate}

While CUDA provides a highly optimized environment for Nvidia hardware, OpenCL was chosen for its cross-platform portability.
This allowed for development and kernel debugging on integrated Intel graphics (laptop) before deploying to high-performance 
Nvidia/AMD discrete GPUs for final benchmarking.

\subsection{Evaluation Methodology}

The project will be evaluated using:

\begin{itemize}
  \item Correctness testing: Unit tests comparing the output of optimised methods against the naive baseline.
  \item Performance Benchmarking: Measuring wall clock time, memory usage, and throughput across various datasets.
  \item Scaling Analysis: Measuring the impact of increasing thread counts to identify the sequential limits of the implementation.
\end{itemize}
